\section{LoRA 在 RL 中的角色：不仅是省算力}

讲者讨论了 LoRA 在 RL 中广泛有效的原因，提出一个重要观点：LoRA 的价值不止是低成本微调，更在于它天然带有 \textbf{正则化} 特性，能抑制策略在 RL 中过度漂移。

在 on-policy/off-policy 混合场景里，采样策略与更新策略差异过大容易导致不稳定。LoRA 通过受限参数更新空间让策略变化更平滑，从而有助于收敛稳定性。

\begin{importantbox}{LoRA for RL 的三重收益}
\begin{itemize}
    \item \textbf{稳定性}：减小策略跳跃，降低训练震荡；
    \item \textbf{效率}：减少显存占用，加快试验周转；
    \item \textbf{可运营性}：支持多实验并行与快速回滚。
\end{itemize}
\end{importantbox}

\begin{knowledgebox}{LoRA 与 ``regularization-first'' 思维}
讲者将 RL 看成更需要约束的训练阶段：模型能力强但目标噪声更大，因此 ``可控更新'' 往往比 ``最大自由度更新'' 更适合工业迭代。
\end{knowledgebox}

需要注意的是，LoRA 并非万能。对于需要大幅能力重构的阶段，full update 仍可能必要。实际工程往往采用分层策略：快速探索阶段偏 LoRA，最终冲刺阶段再局部放开参数。

\begin{warningbox}{常见误区：把 LoRA 当成纯压缩技术}
若只从成本角度看 LoRA，容易忽视其在收敛稳定性上的收益；反之若把 LoRA 绝对化，也可能错失需要大幅参数更新的场景。
\end{warningbox}

\subsection{本章小结}
LoRA 在 Agent RL 中的价值是 ``效率 + 稳定'' 双重收益。它本质上是训练约束工具，而不仅是资源节省手段。

